\setchoices{label=(\Alph*)}

\begin{question}[index=1]
曲线 $\displaystyle y=-x\left[\frac1x\right]+\frac{x^2}{\sqrt{1+x^2}}$ 有多少条渐近线？其中 $[\cdot]$ 表示取整函数。
\begin{choices}[columns=4]
\item $0$
\item $1$
\item $2$
\item $3$
\end{choices}
\end{question}
\QuestionBreak

\begin{question}[index=2]
下列命题正确的是
\begin{choices}[columns=1]
\item 级数 $\displaystyle\sum_{n=1}^{\infty}u_nv_n$ 收敛，则 $\displaystyle\sum_{n=1}^{\infty}(u_n^2+v_n^2)$ 收敛。
\item 正项级数 $\displaystyle\sum_{n=1}^{\infty}u_n$ 收敛，则 $\displaystyle\lim_{n\to\infty}\frac{u_{n+1}}{u_n}\le1$。
\item 级数 $\displaystyle\sum_{n=1}^{\infty}u_n$ 收敛，$\displaystyle\sum_{n=1}^{\infty}v_n$ 发散，则 $\displaystyle\sum_{n=1}^{\infty}u_nv_n$ 一定发散。
\item 级数 $\displaystyle\sum_{n=1}^{\infty}u_n^2$ 收敛，则 $\displaystyle\sum_{n=1}^{\infty}u_nu_{n+1}$ 收敛。
\end{choices}
\end{question}
\QuestionBreak

\begin{question}[index=3]
设曲面 $\Sigma\colon z=\dfrac{x^2+y^2}{2}$（$z\le1$）取下侧，$\Sigma_1$ 是 $\Sigma$ 位于第一卦限的部分，则对坐标的曲面积分
\[
  \iint_{\Sigma} yz\dd y\dd z+xz\dd z\dd x+z^2\dd x\dd y
\]
化为对面积的曲面积分为
\begin{choices}[columns=2]
\item $\displaystyle 4\iint_{\Sigma_1}\frac{z^2-2xyz}{\sqrt{1+2z}}\dd S$
\item $\displaystyle 4\iint_{\Sigma_1}\frac{-z^2}{\sqrt{1+2z}}\dd S$
\item $\displaystyle 2\iint_{\Sigma_1}\frac{2xyz-z^2}{\sqrt{1+2z}}\dd S$
\item $\displaystyle 2\iint_{\Sigma_1}\frac{z^2}{\sqrt{1+2z}}\dd S$
\end{choices}
\end{question}
\QuestionBreak

\begin{question}[index=4]
设 $k,n$ 均为正整数，积分 $\displaystyle I=\int_0^{n\pi}\sqrt{1+\sin 2kx}\dd x$ 的值
\begin{choices}[columns=2]
\item 与 $n,k$ 均无关
\item 与 $n,k$ 均有关
\item 与 $n$ 无关，与 $k$ 有关
\item 与 $n$ 有关，与 $k$ 无关
\end{choices}
\end{question}
\QuestionBreak

\begin{question}[index=5]
设 $\displaystyle f(x)=\begin{vmatrix}
x-1&-2x&1&7\\
3&x-2&0&1\\
4&-5&2x+1&4x+5\\
1&0&3&x+2
\end{vmatrix}$，则 $f^{(3)}(0)=$
\begin{choices}[columns=4]
\item $-6$
\item $0$
\item $-1$
\item $6$
\end{choices}
\end{question}
\QuestionBreak

\begin{question}[index=6]
已知 $\vect\alpha=(a,-1,0)^{\mathrm T}$ 是二次型
\[
 \vect x^{\mathrm T}\vect A\vect x=x_1^2+ax_2^2-x_3^2+2(a-1)x_1x_2+4x_1x_3+4x_2x_3
\]
的矩阵 $\vect A$ 的特征向量。若 $\vect x^{\mathrm T}(\vect A+k\vect E)\vect x=1$ 经可逆变换可化为 $y_1^2-y_2^2=1$，则 $k=$
\begin{choices}[columns=4]
\item $3$
\item $-1$
\item $-3$
\item $1$
\end{choices}
\end{question}
\QuestionBreak

\begin{question}[index=7]
设 $\vect A$ 为 $n$ 阶可逆矩阵，$\lambda=2$ 是 $\vect A^{-1}$ 的特征值，对应的特征向量为 $\vect\alpha$；$\mu=1$ 是 $\vect A^{\mathrm T}$ 的特征值，对应的特征向量为 $\vect\beta$，则必有
\begin{choices}[columns=2]
\item $\vect\beta=2\vect\alpha$
\item $\vect\beta=\dfrac12\vect\alpha$
\item $\vect\alpha$ 与 $\vect\beta$ 正交
\item $\vect\alpha$ 与 $\vect\beta$ 线性无关但不正交
\end{choices}
\end{question}
\QuestionBreak

\begin{question}[index=8]
设随机事件 $A,B,C$，若 $A,C$ 相互独立，$B,C$ 相互独立，又 $A,B$ 互不相容，则 $A\cup B$ 与 $C$ 必定
\begin{choices}[columns=4]
\item 不相互独立
\item 相互独立
\item 互不相容
\item 相容
\end{choices}
\end{question}
\QuestionBreak

\begin{question}[index=9]
设随机变量 $X$ 的分布函数为 $F(x)=\begin{cases}a+b\eu^{-x^2/2},&x>0,\\0,&x\le0,\end{cases}$
$E(X)$ 为 $X$ 的期望，则概率 $P\{X>E(X)\}=$
\begin{choices}[columns=4]
\item $1-\eu^{-\pi/2}$
\item $\eu^{-\pi/2}$
\item $1-\eu^{-\pi/4}$
\item $\eu^{-\pi/4}$
\end{choices}
\end{question}
\QuestionBreak

\begin{question}[index=10]
设总体 $X$ 的概率密度为 $f(x)=\begin{cases}2\lambda x\eu^{-\lambda x^2},&x>0,\\0,&x\le0,\end{cases}$
$\lambda>0$ 未知。对 $X$ 进行独立重复观测 $10$ 次，结果有 $4$ 次观测值大于 $1$，则 $\lambda$ 的极大似然估计值 $\widehat\lambda_L=$
\begin{choices}[columns=4]
\item $\ln 5$
\item $\ln 4$
\item $\ln\dfrac52$
\item $\ln\dfrac53$
\end{choices}
\end{question}
